%#!latexmk -pdfdvi -f main.tex
%#LPR acroread %s

\RequirePackage{plautopatch}
\documentclass[paper=a4,
fontsize=12pt,
reqno,
book,
uplatex,
dvipdfmx]{jlreq}
%
%\jlreqsetup{
%thebibliography_heading={\chapter*{\refname}},
%}

%
%%%%%%%%%%%%%%% Include Package file %%%%%%%%%%%%%%%
\usepackage{amsmath,amssymb,amsthm}
%\usepackage{mathrsfs}
\usepackage{graphicx,color}
\usepackage{wrapfig}
\usepackage{newtxtext}
\usepackage[varg]{newtxmath}
\usepackage[shortlabels]{enumitem}
\usepackage{framed}
%\usepackage{otf}
\usepackage[deluxe]{otf}
\usepackage[unicode,noalphabet]{pxchfon}
\setminchofont{BIZUDMincho-Regular.ttf}
\setgothicfont{BIZUDGothic-Regular.ttf}
\setboldgothicfont{BIZUDGothic-Bold.ttf}

\usepackage{algorithm}
\usepackage{algorithmic}


%\usepackage{input/chapterJP}
%
%%%%%%%%%%%%%%% Paper size settings %%%%%%%%%%%%%%%
%
%
\usepackage{geometry}%
\geometry{left=20mm,right=20mm,top=15mm,bottom=25mm}
\geometry{a4paper}
%


%%%%%%%%%%%%%%% index settings %%%%%%%%%%%%%%%
%
%
\usepackage{makeidx}%
\makeindex

%%%%%%%%%%%%%%% references settings %%%%%%%%%%%%%%%
%
%
\usepackage{amsrefs}
%\usepackage[alphabetic,y2k]{amsrefs}
\renewcommand{\bibname}{参考文献}

%\usepackage[
%backend=biber,
%style=alphabetic,
%]{biblatex}
%\ExecuteBibliographyOptions{sorting=nty,maxnames=8,minnames=1,url=false}
%\addbibresource{references.bib}


%%%%%%%%%%%%%%% mathtools %%%%%%%%%%%%%%%
%
\usepackage{mathtools}
\mathtoolsset{showonlyrefs=true}
%
%
%
%%%%%%%%%%%%%%% hyperref %%%%%%%%%%%%%%%
%
\usepackage{hyperref}
\usepackage{pxjahyper}
\hypersetup{%
 bookmarksnumbered=true,%
 colorlinks=false,%
 setpagesize=false,%
 pdftitle={数学キャリアデザイン},%
 pdfauthor={水野将司},%
 pdfsubject={},%
 pdfkeywords={}
}
%
%
%
%%%%%%%%%%%%%%% Setting of Theorem-like environment %%%%%%%%%%%%%%%
% [Italic font in Theorem-like environment]
%\theoremstyle{definition} % if change the roman font
\theoremstyle{definition} % Change the Roman font
%\newtheorem{definition}{定義}[chapter]
\newtheorem{theorem}{定理}[chapter]
\newtheorem{proposition}[theorem]{命題}
\newtheorem{corollary}[theorem]{系}
\newtheorem{lemma}[theorem]{補題}
\newtheorem{problem}{問題}[chapter]
\newtheorem{example}{例}[chapter]
\newtheorem*{example*}{例}
\newtheorem{definition}{定義}[chapter]
\newtheorem*{remark}{注意}
%
 
%
%%%%%%%%%% Proof %%%%%
\renewcommand{\proofname}{証明}  
% ((Use the environment in proof))
% \begin{proof}
% \end{proof}
% white box in the end of proof
%
%
% make new QED
% \QED{foo}
% we write a black box in the right and output foo after the black box
\newcommand{\QED}[1]{\hfill$\blacksquare$ #1\par}
% 
%
%
%%%%%%%%% 1.1 times Line stretch %%%%%%%%%
% If you do not use Japanese, that command should be commented out.
%
%
%%%%%%% Mathematical symbols %%%%%%%%%%%%%
%
\input{mymacro.tex}
%
%%%%%%%%%%%%%%% Counter setting in the equation's reference %%%%%%%%%%%%%%%
% For amsart, use this counter setting
%
\numberwithin{equation}{chapter}
%
%%% Definition of the \section (not use the \scshape)
%


%%%%%%%%%%%%%%% temporary Package %%%%%%%%%%%%%%%

%\usepackage{showkeys}   % Output of cite and reference
%\usepackage{refcheck}   % check of no uselabel 

%%%%%%%%%%%%%%% title Setting %%%%%%%%%%%%%%%
%
%
\title{数学キャリアデザイン}
\author{水野 将司}
%
%
%
%%%%%%%%%%%%%%% page style setting %%%%%%%%%%%%%%%
%
\allowdisplaybreaks[1]
\DeclareMathOperator{\rot}{rot}
\let\Re\relax
\let\Im\relax
\DeclareMathOperator{\Re}{Re}
\DeclareMathOperator{\Im}{Im}

\usepackage{bm}
\renewcommand{\vec}[1]{\bm{#1}}


%
%%%%%%%%%%%%%%% end preamble %%%%%%%%%%%%%%%
%
\begin{document}
%\maketitle

% \begin{flushright}
%  更新日時 : \the\year.\the\month.\the\day
% \end{flushright}


%#!latexmk -pdfdvi -f main.tex


\chapter{微分方程式の数値計算と散逸則(2026年度)}

散逸則を持つ常微分方程式を例にとり，数値計算法と散逸則との関係をみる．


 \section{常微分方程式}

 次の常微分方程式を考える:
 \begin{equation}
  \tag{ODE}
   \label{eq:2026.ODE}
   \left\{
    \begin{aligned}
     x'(t) &= -f_x(x(t)), & t > 0 \\
     x(0) &= x_0 \in \mathbb{R}
    \end{aligned}
   \right.
 \end{equation}
 ここで，$x = x(t) : [0, \infty) \rightarrow \mathbb{R}$ は未知関数，$f =
 f(x) : \mathbb{R} \rightarrow \mathbb{R}$ は与えられた非負値関数，$x_0
 \in \mathbb{R}$ は与えられた初期値である．

 \eqref{eq:2026.ODE}の解に対して，次の散逸則が成り立つ．

 \begin{proposition}[散逸則(微分型)]
  $x = x(t)$ を \eqref{eq:2026.ODE} の解とすると，すべての$t>0$に対して
  \begin{equation}
   \label{eq:2026.LawEnergyDissipation_DiffForm}
    \frac{d}{dt} f(x(t)) = -|f_x(x(t))|^2
  \end{equation}
  が成り立つ．
 \end{proposition}

 \begin{proof}
  合成関数の微分法と~\eqref{eq:2026.ODE} より，
  \begin{equation}
  \frac{d}{dt} f(x(t))
   =
   f_x(x(t))
   \frac{dx}{dt}(t)
   =
   f_x(x(t))
   \cdot
   (-f_x(x(t)))
   =
   -|f_x(x(t))|^2
  \end{equation}
  が得られる．
 \end{proof}

 \eqref{eq:2026.LawEnergyDissipation_DiffForm}を $0 \le t \le T$ で積分す
 ると，次の積分型の散逸則がわかる:
 \begin{equation}
  \label{eq:2026.LawEnergyDissipation_IntegralForm}
   f(x(T)) = f(x(0)) - \int_{0}^{T} |f_x(x(t))|^2 \, dt
 \end{equation}

 \section{常微分方程式の計算例}

 $f$を具体的に与えて微分方程式~\eqref{eq:2026.ODE}の解を求めてみる．

 \begin{example}
  $f(x) = \frac{1}{2}x^2$とすると， $f_x(x) = x$となるか
  ら，~\eqref{eq:2026.ODE}は
  \begin{equation}
   \left\{
    \begin{aligned}
     x'(t) &= -x(t), & t > 0 \\
     x(0) &= x_0 \in \mathbb{R}
    \end{aligned}
	 \right.
  \end{equation}
  となる．この解は$x(t)=x_0e^{-t}$である．
 \end{example}

 \begin{example}
  $f(x) =  \frac{1}{4}x^4$とすると， $f_x(x) = x^3$となるか
  ら，~\eqref{eq:2026.ODE}は
  \begin{equation}
   \left\{
    \begin{aligned}
     x'(t) &= -(x(t))^3, & t > 0 \\
     x(0) &= x_0 \in \mathbb{R}
    \end{aligned}
	 \right.
  \end{equation}
  となる．これは陽的に解けて
  \begin{equation}
   x(t)
    =
    \begin{cases}
     \frac{1}{\sqrt{2t+x_0^2}} & (x_0>0)\\
     0 & (x_0=0) \\
     -\frac{1}{\sqrt{2t+x_0^2}} & (x_0<0)\\
    \end{cases}
  \end{equation}
  となる．
 \end{example}

 \begin{example}
  $f(x) = 1 + \cos x$とすると，$f_x(x) =-\sin x$となるか
  ら，~\eqref{eq:2026.ODE}は
  \begin{equation}
   \left\{
    \begin{aligned}
     x'(t) &= \sin(x(t)), & t > 0 \\
     x(0) &= x_0 \in \mathbb{R}
    \end{aligned}
	 \right.
  \end{equation}
  となる．これは陽的に解けて
  \begin{equation}
   x(t)
    =
    \begin{cases}
     \frac{1}{\sqrt{2t+x_0^2}} & (x_0>0)\\
     0 & (x_0=0) \\
     -\frac{1}{\sqrt{2t+x_0^2}} & (x_0<0)\\
    \end{cases}
  \end{equation}
  となる．この解は$x(t) = 2 \arctan(\tan(x_0/2) e^t)$である．
 \end{example}

 \begin{example}
  $f(x) = e^{x^2}$とすると，$f_x(x) =2e^{x^2}$となるか
  ら，~\eqref{eq:2026.ODE}は
  \begin{equation}
   \left\{
    \begin{aligned}
     x'(t) &= -2(x(t))e^{(x(t))^2}, & t > 0 \\
     x(0) &= x_0 \in \mathbb{R}
    \end{aligned}
	 \right.
  \end{equation}
  となる．これは陽的に解けないらしい．
 \end{example}


 \section{数値的に求める}

 一般的に~\eqref{eq:2026.ODE}の解は陽的に書けないことの方が多い．そこで，
 微分を差分におきかえて，近似解を探すことを試みる．

 微分係数の定義
 \begin{equation}
  x'(t)
   =
   \lim_{\Delta t \rightarrow 0}
   \frac{x(t+\Delta t) - x(t)}{\Delta t}
 \end{equation}
 において，十分小さい時間刻み幅 $\Delta t > 0$ を固定して，
 \begin{equation}
  \label{eq:2006.Derivative_Approx}
  x'(t) \approx \frac{x(t+\Delta t) - x(t)}{\Delta t}
 \end{equation}
 と近似すると，~\eqref{eq:2026.ODE}は次のように差分化できる:
 \begin{equation}
  \frac{x(t+\Delta t) - x(t)}{\Delta t} = -f_x(x(t))
 \end{equation}
 これを$x(t+\Delta t)$について解いて
 \begin{equation}
  x(t+\Delta t) = x(t) - \Delta t f_x(x(t))
 \end{equation}
 としよう．時間ステップを $t = 0, \Delta t, 2\Delta t, 3\Delta t, \dots$
 と進めると，
 \begin{align}
  x(\Delta t) &= x(0) - \Delta t f_x(x(0)), \\
  x(2\Delta t) &= x(\Delta t) - \Delta t f_x(x(\Delta t)), \\
  x(3\Delta t) &= x(2\Delta t) - \Delta t f_x(x(2\Delta t)),
 \end{align}
 となる。$a_n = x(n\Delta t)$ ($n = 0, 1, 2, \dots$) とおくと，
 次の漸化式を解く問題に変形できる:
\begin{equation}
 \label{eq:2006.FDM}
  \begin{cases}
   a_{n+1} = a_n - \Delta t f_x(a_n) ,\qquad n=0,1,2,\ldots\\
   a_0 = x_0
  \end{cases}
  \tag{RE}
\end{equation}
目的の時間 $T > 0$ における値 $x(T)$ を求めたいときは，$N\Delta t \ge T$
となるようなステップ数 $N$ （すなわち $N \ge \frac{T}{\Delta t}$）を決め
ればよい．

 \section{漸化式をみたす数列を求めるスキーム}

 \eqref{eq:2006.FDM}をみたす数列を求めるスキームの構成は以下の通りとなる．

  \subsection*{与える定数}
  \begin{itemize}
   \item 初期値 $x_0\in\R$
   \item 時間刻み幅 $\Delta t > 0$
   \item 最終時間 $T > 0$
  \end{itemize}
  ここから，総ステップ数$N$を
  \begin{equation}
   N = \text{切り上げ}\left(\frac{T}{\Delta t}\right)
  \end{equation}
  で求める．

  \subsection*{与える関数}
  \begin{itemize}
   \item 非負値関数 $f:\R\rightarrow\R$
  \end{itemize}
  ここから，$f$の導関数$f_x$を計算する．

  \begin{remark}
   $f$は与えられた関数だから，その導関数$f_x$は厳密に計算できる．従って，
   \eqref{eq:2006.Derivative_Approx}のような近似計算は必要ない．しかしな
   がら，導関数を求める計算は「記号計算」と呼ばれるもので，「数値計算」
   とは異なる方法をとる(mathematicaも記号計算ができる)．そこで，pythonの
   プログラミングでは，「記号計算」ができるライブラリであるsympyを使って，
   与えられた$f$から微分を「記号計算」で計算させる．その後，sympyの
   lambdify関数を使って，「記号計算」として定めた関数を「数値計算」に利
   用できるように変換するようにした．
  \end{remark}

  \subsection*{初期化}
  numpy型の配列 $a_0, a_1, a_2, \dots, a_N$を$a_0=a_1=a_2=\ldots
  =a_N=0.0$で初期化する．そのあと，$a_0 = x_0$と置く．

  \subsection*{繰返計算(漸化式を求める)}
  $n = 0, 1, 2, \dots, N-1$ について，以下の代入文を繰り返し実行する:
  \begin{equation}
   a_{n+1} = a_n - \Delta t f_x(a_n).
  \end{equation}


 \section{近似解で散逸則は成り立つか}

 上記で得られた$a_N\approx x(T)$に対して，積分型の散逸
 則~\eqref{eq:2026.LawEnergyDissipation_IntegralForm}が成り立つかどうかを
 調べる．そのために，積分については台形則を用いることにする．すなわち
 \begin{multline}
  \int_0^T|f_x(x(t))|^2\,dt
  \\
   =
   \frac{1}{2}|f_x(x(0))|^2\Delta t
   +
   |f_x(x(\Delta t))|^2\Delta t
   +
   |f_x(x(2\Delta t))|^2\Delta t
   +
   \cdots
   |f_x(x((N-1)\Delta t))|^2\Delta t
   +
   \frac{1}{2}|f_x(x(N\Delta t))|^2\Delta t
  \\
  =
  \left(
  |f_x(x(0))|^2,|f_x(x(\Delta t))|^2,|f_x(x(2\Delta t))|^2,\ldots,|f_x(x(N\Delta t))|^2
  \right)
  \cdot
  \left(
  \frac{1}{2},1,1,\ldots,1,\frac{1}{2}
  \right)
  \Delta t
 \end{multline}
 を用いることにする．~\eqref{eq:2026.LawEnergyDissipation_IntegralForm}が成り立つかどうかをみるために，
\begin{multline}
 \left|
 f(x(T)) -
 \left(
  f(x(0))
  -
  \int_{0}^{T} |f_x(x(t))|^2 \, dt
 \right)
 \right|
 \\
 =
 \left|
 f(a_N) -
 \left(
  f(a_0)
  -
 \left(
  |f_x(a_0)|^2,|f_x(a_1)|^2,|f_x(a_2)|^2,\ldots,|f_x(a_N)|^2
  \right)
  \cdot
  \left(
  \frac{1}{2},1,1,\ldots,1,\frac{1}{2}
  \right)
  \Delta t
 \right)
 \right|
\end{multline}
を計算させてみよう．

 \section{アルゴリズムの詳細と実装}
 
 % 必要に応じて、日本語環境用の設定
 \renewcommand{\algorithmicrequire}{\textbf{Input:}}
 \renewcommand{\algorithmicensure}{\textbf{Output:}}
 
 本稿で用いた，常微分方程式~\eqref{eq:2026.ODE}の差分法(前進オイラー法)に
 よる数値シミュレーション，およびエネルギー関数の時間変化(台形公式による
 積分)を検証するための具体的な計算手順を以下に示す．

  \subsection{計算アルゴリズム}
  
  時間ステップ幅を $\Delta t$，最終時刻を $T$ とし，反復回数を $N =
  \lceil T / \Delta t \rceil$ とする．初期値 $x_0$ から出発し，関数
  $f(x)$ に対する速度場 $v(x) = - \frac{df}{dx}$ を用いて状態 $a_n
  \approx x(n\Delta t)$ を更新する．また，エネルギーの保存(減少)誤差を台
  形公式による数値積分を用いて評価する．
  
  \begin{algorithm}[H]
   \caption{前進オイラー法による軌道計算とエネルギー変化の検証}
   \label{alg:forward_euler}
   \begin{algorithmic}[1]
    \REQUIRE 初期値 $x_0$, 時間刻み幅 $\Delta t$, 最終時刻 $T$, 評価対象の関数 $f(x)$
    \ENSURE 状態の履歴 $\{a_n\}_{n=0}^N$, エネルギー誤差 $\mathcal{E}$
    \STATE $N \leftarrow \lceil T / \Delta t \rceil$ \COMMENT{総反復回数の決定}
    \STATE $v(x) \leftarrow - \frac{df}{dx}$ \COMMENT{速度関数（導関数）の定義}
    \STATE 配列 $a$ をサイズ $N+1$ で初期化
    \STATE $a_0 \leftarrow x_0$ \COMMENT{初期条件の設定}
    \FOR{$n = 0$ \TO $N-1$}
    \STATE $a_{n+1} \leftarrow a_n + \Delta t \cdot v(a_n)$ \COMMENT{前進オイラー法による更新}
    \ENDFOR
    \STATE 各時刻 $t_n = n \cdot \Delta t$ ($n=0, \dots, N$) に対する時間配列 $\{t_n\}$ を生成
    \STATE $y_n \leftarrow |v(a_n)|^2$ ($n=0, \dots, N$) \COMMENT{速度の2乗の計算}
    \STATE $I \leftarrow \Delta t \left( \frac{1}{2}y_0 + \sum_{n=1}^{N-1} y_n + \frac{1}{2}y_N \right)$ \COMMENT{台形公式による $\int_0^T |v(x(t))|^2 dt$ の近似計算}
    \STATE $\mathcal{E} \leftarrow \left| f(a_N) - f(0) + I \right|$ \COMMENT{エネルギー誤差の計算}
\RETURN $\{a_n\}_{n=0}^N$, $\mathcal{E}$
   \end{algorithmic}
  \end{algorithm}

  \subsection{プログラムの構成と処理の流れ}

  上記のアルゴリズムを以下の3つのフェーズで実装した．
  
  \begin{enumerate}
   \item \textbf{シンボリック計算（SymPyによる自動微分）}:
	 %
	 評価したい関数 $f(x) = \exp(x^2)$ を定義し，その導関数の $-1$
	 倍である $v(x) = -2x\exp(x^2)$ を自動微分により算出する．その後，
	 \texttt{lambdify} を用いて高速な数値計算（NumPy）が可能な形式へ
	 と変換する．
	 
   \item \textbf{時間発展のシミュレーション（NumPy）}:
	 %
	 前進オイラー法（Forward Euler Method）に基づき，漸化式 $a_{n+1}
	 = a_n + \Delta t \cdot v(a_n)$ をループ処理によって繰り返し計算
	 し，状態変数 $x(t)$ の時系列データを取得する．
    
    \item \textbf{エネルギー解析と可視化（Matplotlib）}:
	  %
	  得られた軌道から，各時刻におけるエネルギー $f(x(t))$ の減衰挙
	 動をプロットする．さらに，エネルギーの減少量と散逸項の積分値(台
	 形公式で計算)を比較し，理論解からの数値的な乖離を表す「エネルギー
	 誤差」を出力する．
  \end{enumerate}

\section{コメント}

常微分方程式の数値計算については，
\cites{BC08225057,BA32760992,BB20100756}を参考にせよ．本稿でのソースコー
ドはpythonを用いた．ソースコード中のライブラリや関数については，
\cite{BB30665488}を参考にせよ．


% bibtex Settings
\bibliography{references}


\printindex

\end{document}

% Local Variables:
% mode: yatex
% TeX-master: t
% End:

